%% ma: L12-B08-0140
%% lop: 12
%% chuong: 2
%% bai: 8
%% dang: 12.8.27
%% loai: TN4
%% muc: VD
%% dap_an: "A"
%% nguon: "CĐ2.3.pdf giáo viên gửi 27/10/2026 (M07, Lê Minh Khôi – Nguyễn Anh Khôi, Chương II Chủ đề 3 Biểu thức tọa độ của các phép toán vectơ), Câu 56"
%% kien_thuc: "Bài 8 (trọng tâm tứ diện, trọng tâm tam giác, độ dài), Bài 6 (tính chất trọng tâm tứ diện)"
%% trang_thai: chinh-thuc
\begin{ex}
Trong không gian $Oxyz$, cho tứ diện $ABCD$. Gọi $I$, $G$ lần lượt là trọng tâm của tứ diện $ABCD$ và tam giác $BCD$. Biết rằng tọa độ $A(1;2;-1)$ và $G(4;-5;6)$. Khoảng cách $OI$ bằng bao nhiêu?
\choice{\True $OI=\dfrac{\sqrt{627}}4$}{$OI=6$}{$OI=\dfrac{\sqrt{627}}2$}{$OI=\sqrt{627}$}
\loigiai{$\overrightarrow{AI}=\dfrac34\overrightarrow{AG}=\dfrac34(3;-7;7)\Rightarrow I\left(\dfrac{13}4;-\dfrac{13}4;\dfrac{17}4\right)$.

$OI=\dfrac{\sqrt{169+169+289}}4=\dfrac{\sqrt{627}}4$.}
\end{ex}
